\subsection*{Residual structure and the scope of the classification}
The codimension formula leads to a description beyond the stable
component theorem retained in the appendices. Theorem~\ref{thm:all-dimension}
removes its dimension restriction: in the original range $4\le c<k$
there are exactly five nonendpoint equality incidences. In each the
plane contains the entire radical and is maximal isotropic in the
quotient. Orientation monodromy, rather than a choice of one connected
orthogonal family, determines their irreducibility.

The common mechanism is Theorem~\ref{thm:residual-tangent}. For an
annihilator $\ell$, put $W_\nu=U_\nu\cap\operatorname{rad}H_\ell$.
The obstruction to solving the infinitesimal plane equations is exactly
the restriction of $\dot\ell$ to $\sum_\nu W_\nu^2$.
This statement holds at every annihilator rank. At rank two it turns
singularity of a corank-one point into failure of a smaller multiplication
map. At the wall a general point of the residual determinant divisor
is an ordinary double crossing between the nondegenerate and signed
components (Theorem~\ref{thm:wall-node}), with a logarithmic conditioning
law (Corollary~\ref{cor:node-tail}).

The reduced rank-two incidence has a normalization indexed by sign
patterns. Its intersections with a common annihilator have a complete
relative-Grassmannian description, while its full isotropy fibres carry
an explicit square-zero nilradical
(Theorem~\ref{thm:signed-normalization} and
Proposition~\ref{prop:rank-two-fibre}). These are different scheme
statements. A reduced global scheme need not have reduced fibres.
For higher symmetric powers the analogous labels form dihedral phase
orbits, rather than a simple power of the number of hyperplanes
(Theorems~\ref{thm:higher-phase} and \ref{thm:higher-residual}).

The all-dimension component theorem is proved from the codimension
formula and the residual tangent sequence; it does not rely on the
stable theorem retained in the appendices. The full boundary,
count-comparison and finite-sample arguments are also retained there,
so that their statistical assumptions do not enter the geometric proof.

Our use of secant obstructions is situated in the classical failure-locus
literature, including Ballico's higher-order failure-locus paper
\cite{Ballico93}, the complete-series work of Green--Lazarsfeld
\cite{GL86}, and the failure-cycle study \cite{Ballico96}.
The statements proved here specify the varying subseries, the absence
of cross-contact products, the exceptional incidences, and the local
scheme strata. The accompanying literature audit records which primary
texts were checked and which theorem-level comparisons remain
unverified. For \cite{Ballico93}, the first page has now been inspected in addition
to the bibliographic record; this is not a completed theorem-by-theorem
comparison. Section~\ref{sec:prior-comparison} specifies this distinction.

For clarity, the Hankel shape identification used below is an equality
of ideals in the same polynomial ring, not just a dimension analogy.
Writing $S=\C[z_0,\ldots,z_{2k-2}]$ and
$H_{u,v}=(z_{i+j})_{0\le i<u,\,0\le j<v}$ for $u+v=2k$,
the classical identity recalled in \cite[Section 1]{CMSV} is
\[
 I_{r+1}(H_{k,k})=I_{r+1}(H_{r+1,\,2k-r-1})\subset S
       \qquad(1\le r<k).
\]
The latter is a maximal-minor Hankel ideal with quotient of affine
dimension $2r$. Projectivization gives $2r-1$; removing the smaller
rank locus gives the exact-rank stratum used in the incidence proof.
